% In this file you should put all LaTeX macros and settings to be used both by
% the pdf version and the web version.
% This should be most of your macros.

% The theorem-like environments defined below are those that appear by default
% in the dependency graph. See the README of leanblueprint if you need help to
% customize this. 
% The configuration below use the theorem counter for all those environments
% (this is what the [theorem] arguments mean) and never resets it.
% If you want for instance to number them within chapters then you can add
% [chapter] at the end of the next line.
% Final style grid (user decree, M4 addenda #1+#3): ITALIC body for every
% content/claim environment --- theorem, proposition, lemma, corollary,
% conjecture, definition, abbreviation (amsthm's "plain" style: bold/bold
% heading, italic body) --- and UPRIGHT body for every discourse
% environment --- proof (amsthm's built-in default, unaffected by
% \theoremstyle), example, remark. "plain" is amsthm's own default before
% any \theoremstyle call, but it is set explicitly here for clarity.
\theoremstyle{plain}
\newtheorem{theorem}{Theorem}
\newtheorem{proposition}[theorem]{Proposition}
\newtheorem{lemma}[theorem]{Lemma}
\newtheorem{corollary}[theorem]{Corollary}
% `conjecture` (M4 addendum #2): a claim, not yet a theorem --- same
% italic style and numbering family as the rest of this group, but NEVER
% followed by a \begin{proof}: a conjecture environment never carries a
% proof, and never carries a proof-side \leanok (there is nothing to mark
% formalized). Its one current use is \label{cnj:chain-bound} in
% content.tex (reclassified from the former corollary
% \label{cor:chain-bound}, per that addendum).
\newtheorem{conjecture}[theorem]{Conjecture}
\newtheorem{definition}[theorem]{Definition}

% `example` and `remark` use amsthm's "definition" style (bold heading,
% UPRIGHT/roman body) instead: discourse, not a claim or a definition ---
% user decree (M4 addendum #1). `remark` is not one of the numbered
% theorem-like environments above; it lives here alongside them so every
% environment declaration for this document is in one place (moved from
% content.tex, M1).
\theoremstyle{definition}
\newtheorem{example}[theorem]{Example}
\newtheorem{remark}[theorem]{Remark}

% `abbreviation` is an unnumbered, bold run-in "Abbreviation." block,
% ITALIC body like the rest of the content/claim group above (M4 addendum
% #3; amsthm's "plain" style, starred form): reserved for a block that
% introduces a *name* for existing content (a Lean `abbrev` twin) ---
% see the Blueprint-to-Lean correspondence paragraph in content.tex
% (M1), which also states this environment's invariant. Declared here
% but not yet used by any item in this document: its first uses are a
% later ticket's job (re-homing `Tens`/`LogTens` into it). If a future
% plasTeX version chokes on the starred form, fall back to
% \newtheorem{abbreviation}[theorem]{Abbreviation} here --- and note
% that a starred amsthm environment carries no counter of its own, so
% any \label placed inside one needs the numbered fallback before it
% can be \ref'd correctly (observed directly while drafting this
% ticket: plasTeX renders such a \ref as literal "??", and xelatex
% resolves it to the nonsensical enclosing section number).
\theoremstyle{plain}
\newtheorem*{abbreviation}{Abbreviation}

% ----------------------------------------------------------------------
% Shorthands for the NeSyCat library content (blueprint/src/content.tex).
% Kept deliberately at amsmath/amssymb level so plasTeX digests them.
% ----------------------------------------------------------------------

% ===== SHARED-MACRO IMPORT (single source of truth: NeSyCat.Logics/macros.sty) =====
% User decree: symbols shared with the NeSyCat papers (\seq, \seqop,
% \bind, \sem, and the domain-signature vocabulary
% \Rel/\Var/\Fun/\Conn/\inn/\out) are defined ONCE, centrally, in
% NeSyCat.Logics/macros.sty, and imported here rather than re-invented
% locally. The two engines that process this file tolerate the import
% differently (verified directly while drafting this ticket):
%   - real xelatex (pdf build, print.tex): `\usepackage{<path>/macros}`
%     (path without the .sty extension) loads the file cleanly through
%     the normal LaTeX package machinery.
%   - plasTeX (web build, web.tex): the very same `\usepackage` call
%     crashes hard --- `ModuleNotFoundError: No module named
%     'leanblueprint.Packages.'` --- because plasTeX intercepts
%     \usepackage to look for a Python plugin module matching the
%     package name, and a relative path containing a dot
%     (`NeSyCat.Logics`) breaks its module-name parsing. A bare
%     `\input{<path>/macros.sty}` does not crash plasTeX, but \seq/\sem
%     depend on stmaryrd's \fatsemi/\llbracket/\rrbracket, which the web
%     build's plain MathJax config (no stmaryrd extension loaded) is not
%     guaranteed to render. The route below is the sanctioned one: a
%     real import on the pdf path, and on the web path a MINIMAL set of
%     local `\providecommand`s (one-liners, each marked "WEB FALLBACK")
%     whose bodies mirror macros.sty's --- exactly, for \seqop/\bind
%     (kerned with the plain \mkern primitive, MathJax-safe) and via a
%     stmaryrd-free approximation for \seq/\sem (as this file already
%     did, pre-centralization).
\makeatletter
\@ifundefined{XeTeXversion}{%
  % ---- WEB (plasTeX/MathJax) branch: minimal local fallbacks ----
  % \parr: literal Unicode PARR character (U+214B); MathJax/plasTeX
  % already has the glyph on this path (no stmaryrd dependency) ---
  % same disclosed route as before centralization.
  \newcommand{\parr}{\mathbin{⅋}}
  \providecommand{\Rel}{\mathrm{Rel}}                        % WEB FALLBACK (rendering only, defs mirror macros.sty)
  \providecommand{\Var}{\mathrm{Var}}                        % WEB FALLBACK (rendering only, defs mirror macros.sty)
  \providecommand{\Fun}{\mathrm{Fun}}                        % WEB FALLBACK (rendering only, defs mirror macros.sty)
  \providecommand{\Conn}{\mathrm{Conn}}                      % WEB FALLBACK (rendering only, defs mirror macros.sty)
  \providecommand{\inn}{\mathrm{in}}                         % WEB FALLBACK (rendering only, defs mirror macros.sty)
  \providecommand{\out}{\mathrm{out}}                        % WEB FALLBACK (rendering only, defs mirror macros.sty)
  \providecommand{\sem}[1]{[\![#1]\!]}                       % WEB FALLBACK (rendering only, defs mirror macros.sty; stmaryrd-free stand-in for \llbracket/\rrbracket)
  \providecommand{\seq}{\mathbin{;}}                         % WEB FALLBACK (rendering only, defs mirror macros.sty; stmaryrd-free stand-in for \fatsemi)
  \providecommand{\seqop}{\mathrel{>\mkern-4mu>\mkern-4mu>}} % WEB FALLBACK (rendering only, defs mirror macros.sty exactly)
  \providecommand{\bind}{\mathbin{>\mkern-5mu>\mkern-1mu=}}  % WEB FALLBACK (rendering only, defs mirror macros.sty exactly)
}{%
  % ---- PDF (real xelatex) branch: real central import ----
  % stmaryrd supplies \fatsemi and \llbracket/\rrbracket, which \seq and
  % \sem depend on in macros.sty. cmll supplies the real linear-logic
  % \parr glyph that new.tex itself loads it for (`\let\parr=\invamp`);
  % load it if this TeX install has it (confirmed available: its fonts
  % loaded cleanly compiling new.tex directly for this ticket), else
  % degrade to the previous rotated-ampersand approximation.
  \usepackage{stmaryrd}
  \IfFileExists{cmll.sty}{\usepackage{cmll}}{%
    \usepackage{graphicx}
    \newcommand{\parr}{\mathbin{\rotatebox[origin=c]{180}{$\&$}}}
  }
  \usepackage{../../../NeSyCat.Logics/macros}
}
\makeatother
% ===== END SHARED-MACRO IMPORT =====

% ===== SYMBOL INVENTORY (mirror: NeSyCat/Notation.lean, future) =====
% Purely organizational, no rendering effect: this registry mirrors,
% macro-by-macro, what a future `NeSyCat/Notation.lean` would declare as
% Lean `notation`/`scoped` twins; that file does not exist yet, and no
% `\lean{}` mark attaches to any macro here (symbols are not citable
% objects --- see the Blueprint-to-Lean correspondence paragraph in
% content.tex). Entries whose definition now lives centrally carry a
% "def: central" pointer instead of a local \newcommand (see the import
% above); \Dom is the one deliberate exception, disclosed below.
%   \seq    Lean: scoped notation (planned), mirrors mathlib's `≫`      -- def: central (NeSyCat.Logics/macros.sty)
%   \seqop  Lean: scoped infixl (planned) "⋗", feed-forward composite   -- def: central (NeSyCat.Logics/macros.sty)
%   \bind   Lean: scoped infixl (planned) "≫=", Kleisli bind            -- def: central (NeSyCat.Logics/macros.sty)
%   \sem    Lean: scoped notation (planned) "⟦_⟧"                       -- def: central (NeSyCat.Logics/macros.sty)
%   \Rel, \Var, \Fun, \Conn, \inn, \out (domain-signature vocabulary)   -- def: central (NeSyCat.Logics/macros.sty)
% Truth-value-level (logical) monoid connectives and their units, Table
% 3: \parr is or-like (unit $\dot0$), \AndC is and-like (unit $\dot1$).
% Distinct from the categorical/semiring \oplus,\otimes used for
% weight-semiring operations inside $\MS{S}$ (Chapter 1) --- see
% Remark~\ref{rem:connective-conventions} in content.tex, right after
% the Sources-and-provenance paragraph. \parr's engine-guarded route is
% defined above, alongside the shared-macro import; none of \AndC,
% \dzero, \done, \impc, \negc, \actop below exist in macros.sty, so they
% stay locally defined here (def: local).
\newcommand{\AndC}{\mathbin{\&}}                % Lean: scoped infixl " & "                -- def: local
\newcommand{\dzero}{\dot 0}                     % Lean: scoped notation "0̇"                -- def: local
\newcommand{\done}{\dot 1}                      % Lean: scoped notation "1̇"                -- def: local
\newcommand{\impc}[1]{\mathbin{\to_{#1}}}       % Lean: scoped notation:25 a " →[" c "] " b -- def: local
\newcommand{\negc}[1]{\neg_{#1}}                % Lean: scoped prefix:75 "¬[" c "] "        -- def: local
% NeSyCat Theory v2 (M2): the grammatical layer's actegory action. Not
% formalized yet (no Lean file exists), so the twin comment records only
% the planned shape, following the same convention as the rest of this
% inventory.
\newcommand{\actop}{\mathbin{\bullet}}          % Lean: scoped infixr (planned) " • " --- actegory action (Definition~def:categorical-interpretation) -- def: local
% ===== END SYMBOL INVENTORY =====

% Monadic programming notation.
\newcommand{\Ret}{\mathrm{return}}
\newcommand{\Do}{\mathbf{do}}

% Semiring-monad family.
\newcommand{\MS}[1]{\mathcal M_{#1}}
\newcommand{\Idmon}{\mathrm{Id}}
\newcommand{\Dmon}{\mathcal D}
\newcommand{\Tmon}{\mathcal T}
\newcommand{\LTmon}{\log\mathcal T}
\newcommand{\Bmon}{\mathcal B}
\newcommand{\Bidx}{\underline B}
\newcommand{\dstL}{\mathrm{dst}_L}
\newcommand{\dstR}{\mathrm{dst}_R}

% Semiring carriers.
\newcommand{\BoolS}{\mathbb B}
\newcommand{\MassS}{\mathbb R_{\ge 0}}
\newcommand{\LogS}{\mathbb R \cup \{-\infty\}}
\newcommand{\ProbS}{[0,1]}

% Bridge maps and normalization (Chapter 6).
\newcommand{\dec}{\mathrm{dec}}
\newcommand{\enc}{\mathrm{enc}}
\newcommand{\nrm}{\mathrm{nrm}}
\newcommand{\tilt}{\mathrm{tilt}}
\newcommand{\lse}{\operatorname{lse}}
\DeclareMathOperator{\supp}{supp}

% Syntax and semantics (Chapter 4).
% \Dom is intentionally NOT centralized: NeSyCat.Logics/macros.sty
% already uses the name \Dom for an unrelated unary macro
% (\Dom[1]{\mathcal I(s_{#1})}, a different paper's indexed-state-space
% notation), whereas this blueprint's \Dom is the nullary domain-symbol
% set `\mathrm{Dom}` used throughout content.tex's domain-signature
% material (Definition~def:domain-signature onward). Importing the
% central file's \Dom in its place would silently turn every bare
% `$\Dom$` in content.tex into a "Missing argument" error, so this one
% symbol stays a deliberate local override (noticed during the M3
% macro-centralization ticket, not resolved by it: resolving it means
% either renaming content.tex's usage or renaming macros.sty's unary
% \Dom, both out of this ticket's scope). On the pdf path the import
% above already binds \Dom (to the unary macro), so \let it back to
% \relax first --- \newcommand then (re)defines it uniformly on both
% engines without an "already defined" clash.
\let\Dom\relax
\newcommand{\Dom}{\mathrm{Dom}}
\newcommand{\Quan}{\mathrm{Quan}}